\documentclass[10pt,a4paper]{amsart}
\usepackage{amsmath,amssymb,amsthm,graphicx}
% Retained preamble for the Stacks Project chapter {CH} (shared project preamble).
% The chapter file pinned-original-source.tex only carries \input{{preamble}} on its first line,
% so no preamble text lives inside the chapter itself.  The Stacks Project uses one single shared
% preamble for every chapter: the snapshot stages it three times (sources/stacks-chow/preamble.tex,
% sources/stacks-etale-cohomology/preamble.tex and focused-seeds/AGP/sources/stacks-snapshot-20261011/
% preamble.tex) and those copies are identical in size.  The block below is copied mechanically,
% line for line, from sources/stacks-chow/preamble.tex lines 14-18 and 154-217: the Xy-pic setup and
% the native theorem declarations together with the mathematical macros.  Nothing was added, removed
% or redefined; the wrapper class and the amsmath/amssymb/amsthm/graphicx base packages are supplied
% by the assembler.
% For commutative diagrams we use Xy-pic
% Chapter: stacks-spaces-more-morphisms.tex -- More on Morphisms of Spaces.
\usepackage[all]{xy}

% We use 2cell for 2-commutative diagrams.
\xyoption{2cell}

% Theorem environments.
%
\theoremstyle{plain}
\newtheorem{theorem}[subsection]{Theorem}
\newtheorem{proposition}[subsection]{Proposition}
\newtheorem{lemma}[subsection]{Lemma}

\theoremstyle{definition}
\newtheorem{definition}[subsection]{Definition}
\newtheorem{example}[subsection]{Example}
\newtheorem{exercise}[subsection]{Exercise}
\newtheorem{situation}[subsection]{Situation}

\theoremstyle{remark}
\newtheorem{remark}[subsection]{Remark}
\newtheorem{remarks}[subsection]{Remarks}

\numberwithin{equation}{subsection}

% Macros
%
\def\lim{\mathop{\mathrm{lim}}\nolimits}
\def\colim{\mathop{\mathrm{colim}}\nolimits}
\def\Spec{\mathop{\mathrm{Spec}}}
\def\Hom{\mathop{\mathrm{Hom}}\nolimits}
\def\Ext{\mathop{\mathrm{Ext}}\nolimits}
\def\SheafHom{\mathop{\mathcal{H}\!\mathit{om}}\nolimits}
\def\SheafExt{\mathop{\mathcal{E}\!\mathit{xt}}\nolimits}
\def\Sch{\mathit{Sch}}
\def\Mor{\mathop{\mathrm{Mor}}\nolimits}
\def\Ob{\mathop{\mathrm{Ob}}\nolimits}
\def\Sh{\mathop{\mathit{Sh}}\nolimits}
\def\NL{\mathop{N\!L}\nolimits}
\def\CH{\mathop{\mathrm{CH}}\nolimits}
\def\proetale{{pro\text{-}\acute{e}tale}}
\def\etale{{\acute{e}tale}}
\def\QCoh{\mathit{QCoh}}
\def\Ker{\mathop{\mathrm{Ker}}}
\def\Im{\mathop{\mathrm{Im}}}
\def\Coker{\mathop{\mathrm{Coker}}}
\def\Coim{\mathop{\mathrm{Coim}}}

% Boxtimes
%
\DeclareMathSymbol{\boxtimes}{\mathbin}{AMSa}{"02}

%
% Macros for moduli stacks/spaces
%
\def\QCohstack{\mathcal{QC}\!\mathit{oh}}
\def\Cohstack{\mathcal{C}\!\mathit{oh}}
\def\Spacesstack{\mathcal{S}\!\mathit{paces}}
\def\Quotfunctor{\mathrm{Quot}}
\def\Hilbfunctor{\mathrm{Hilb}}
\def\Curvesstack{\mathcal{C}\!\mathit{urves}}
\def\Polarizedstack{\mathcal{P}\!\mathit{olarized}}
\def\Complexesstack{\mathcal{C}\!\mathit{omplexes}}
% \Pic is the operator that assigns to X its picard group, usage \Pic(X)
% \Picardstack_{X/B} denotes the Picard stack of X over B
% \Picardfunctor_{X/B} denotes the Picard functor of X over B
\def\Pic{\mathop{\mathrm{Pic}}\nolimits}
\def\Picardstack{\mathcal{P}\!\mathit{ic}}
\def\Picardfunctor{\mathrm{Pic}}
\def\Deformationcategory{\mathcal{D}\!\mathit{ef}}
\usepackage{cleveref}
\crefname{originaltheorem}{original theorem}{original theorems}
\Crefname{originaltheorem}{Original theorem}{Original theorems}
\makeatletter
\newlabel{etale-cohomology-theorem-topological-invariance}{{etale-cohomology-theorem-topological-invariance}{source}{}{sourcepointer-1}{}}
\newlabel{etale-cohomology-theorem-topological-invariance@cref}{{[originaltheorem][0][]etale-cohomology-theorem-topological-invariance}{[1][0][]source}{}{}{}}
\makeatother
\crefname{originaltheorem}{original theorem}{original theorems}
\Crefname{originaltheorem}{Original theorem}{Original theorems}
\makeatletter
\newlabel{spaces-descent-lemma-fpqc-universal-effective-epimorphisms}{{spaces-descent-lemma-fpqc-universal-effective-epimorphisms}{source}{}{sourcepointer-2}{}}
\newlabel{spaces-descent-lemma-fpqc-universal-effective-epimorphisms@cref}{{[originaltheorem][0][]spaces-descent-lemma-fpqc-universal-effective-epimorphisms}{[1][0][]source}{}{}{}}
\makeatother
\crefname{originaltheorem}{original theorem}{original theorems}
\Crefname{originaltheorem}{Original theorem}{Original theorems}
\makeatletter
\newlabel{spaces-morphisms-lemma-etale-universally-injective-open}{{spaces-morphisms-lemma-etale-universally-injective-open}{source}{}{sourcepointer-3}{}}
\newlabel{spaces-morphisms-lemma-etale-universally-injective-open@cref}{{[originaltheorem][0][]spaces-morphisms-lemma-etale-universally-injective-open}{[1][0][]source}{}{}{}}
\makeatother
\crefname{originaltheorem}{original theorem}{original theorems}
\Crefname{originaltheorem}{Original theorem}{Original theorems}
\makeatletter
\newlabel{spaces-morphisms-section-universal-homeomorphisms}{{spaces-morphisms-section-universal-homeomorphisms}{source}{}{sourcepointer-4}{}}
\newlabel{spaces-morphisms-section-universal-homeomorphisms@cref}{{[originaltheorem][0][]spaces-morphisms-section-universal-homeomorphisms}{[1][0][]source}{}{}{}}
\makeatother
\crefname{originaltheorem}{original theorem}{original theorems}
\Crefname{originaltheorem}{Original theorem}{Original theorems}
\makeatletter
\newlabel{spaces-theorem-presentation}{{spaces-theorem-presentation}{source}{}{sourcepointer-5}{}}
\newlabel{spaces-theorem-presentation@cref}{{[originaltheorem][0][]spaces-theorem-presentation}{[1][0][]source}{}{}{}}
\makeatother
\title{An integral universally injective surjective morphism of algebraic spaces induces an equivalence of small etale topoi}
\author{The Stacks Project Authors}
\date{stacks-project:spaces-more-morphisms.tex:local-snapshot-20261011; GNU FDL 1.2 or later}
\begin{document}
\maketitle
\noindent\textit{Editorial note.} Statement span is lines 1315-1326; proof span is lines 1328-1412, the complete author proof of 85 lines.\par\medskip
\noindent\textit{Editorial note.} The proof contains no xymatrix display; the retained preamble is kept for the shared macro definitions.\par\medskip
\noindent\textit{Editorial note.} The three cited cross-chapter results are prerequisites and are supplied as non-load-bearing navigation stubs.\par\medskip
\noindent\textit{Editorial note.} extra\_packages is empty on purpose: the retained preamble already loads xy with options and a second declaration clashes.\par\medskip
\noindent\textit{Allowed standard prerequisites (editorial).}\begin{itemize}
\item Explicit prerequisite boundary. Integral morphisms of algebraic spaces, universal homeomorphisms, the small etale site of an algebraic space, and descent of etale objects along a representable universal homeomorphism. The faithfulness, fullness and essential surjectivity arguments are retained.
\end{itemize}\medskip
\ifdefined\hypertarget\hypertarget{sourcepointer-1}{}\fi
\noindent\textit{Source reference.} Original theorem etale-cohomology-theorem-topological-invariance: sources/stacks-spaces-more-morphisms/pinned-original-source.tex, source lines 1412--1412. Referenced by the retained statement or proof but not part of the collected spans; supplied only as a navigation pointer. The argument does not depend on its content.\par\medskip
\ifdefined\hypertarget\hypertarget{sourcepointer-2}{}\fi
\noindent\textit{Source reference.} Original theorem spaces-descent-lemma-fpqc-universal-effective-epimorphisms: sources/stacks-spaces-more-morphisms/pinned-original-source.tex, source lines 1412--1412. Referenced by the retained statement or proof but not part of the collected spans; supplied only as a navigation pointer. The argument does not depend on its content.\par\medskip
\ifdefined\hypertarget\hypertarget{sourcepointer-3}{}\fi
\noindent\textit{Source reference.} Original theorem spaces-morphisms-lemma-etale-universally-injective-open: sources/stacks-spaces-more-morphisms/pinned-original-source.tex, source lines 1412--1412. Referenced by the retained statement or proof but not part of the collected spans; supplied only as a navigation pointer. The argument does not depend on its content.\par\medskip
\ifdefined\hypertarget\hypertarget{sourcepointer-4}{}\fi
\noindent\textit{Source reference.} Original theorem spaces-morphisms-section-universal-homeomorphisms: sources/stacks-spaces-more-morphisms/pinned-original-source.tex, source lines 1412--1412. Referenced by the retained statement or proof but not part of the collected spans; supplied only as a navigation pointer. The argument does not depend on its content.\par\medskip
\ifdefined\hypertarget\hypertarget{sourcepointer-5}{}\fi
\noindent\textit{Source reference.} Original theorem spaces-theorem-presentation: sources/stacks-spaces-more-morphisms/pinned-original-source.tex, source lines 1412--1412. Referenced by the retained statement or proof but not part of the collected spans; supplied only as a navigation pointer. The argument does not depend on its content.\par\medskip
\begin{theorem}
\label{theorem-topological-invariance}
Let $S$ be a scheme.
Let $f : X \to Y$ be a morphism of algebraic spaces over $S$.
Assume $f$ is integral, universally injective and surjective.
The functor
$$
V \longmapsto V_X = X \times_Y V
$$
defines an equivalence of categories
$Y_{spaces, \etale} \to X_{spaces, \etale}$.
\end{theorem}

\begin{proof}
The morphism $f$ is representable and a universal homeomorphism, see
Morphisms of Spaces,
Section \ref{spaces-morphisms-section-universal-homeomorphisms}.

\medskip\noindent
We first prove that the functor is faithful.
Suppose that $V', V$ are objects of $Y_{spaces, \etale}$ and
that $a, b : V' \to V$ are distinct morphisms over $Y$.
Since $V', V$ are \'etale over $Y$ the equalizer
$$
E =  V' \times_{(a, b), V \times_Y V, \Delta_{V/Y}} V
$$
of $a, b$ is \'etale over $Y$ also. Hence $E \to V'$ is an \'etale monomorphism
(i.e., an open immersion) which is an isomorphism if and only if it is
surjective. Since $X \to Y$ is a universal homeomorphism we see that this
is the case if and only if $E_X = V'_X$, i.e., if and only if $a_X = b_X$.

\medskip\noindent
Next, we prove that the functor is fully faithful.
Suppose that $V', V$ are objects of $Y_{spaces, \etale}$ and
that $c : V'_X \to V_X$ is a morphism over $X$. We want to construct
a morphism $a : V' \to V$ over $Y$ such that $a_X = c$.
Let $a' : V'' \to V'$ be a surjective \'etale morphism such that $V''$ is
a separated algebraic space. If we can construct a morphism
$a'' : V'' \to V$ such that $a''_X = c \circ a'_X$, then the two compositions
$$
V'' \times_{V'} V'' \xrightarrow{\text{pr}_i} V'' \xrightarrow{a''} V
$$
will be equal by the faithfulness of the functor proved in the first
paragraph. Hence $a''$ will factor through a unique morphism
$a : V' \to V$ as $V'$ is (as a sheaf) the quotient of $V''$ by
the equivalence relation $V'' \times_{V'} V''$. Hence we may assume that
$V'$ is separated. In this case the graph
$$
\Gamma_c \subset (V' \times_Y V)_X
$$
is open and closed (details omitted). Since $X \to Y$ is a universal
homeomorphism, there exists an open and closed subspace
$\Gamma \subset V' \times_Y V$ such that $\Gamma_X = \Gamma_c$.
The projection $\Gamma \to V'$ is an \'etale morphism whose base
change to $X$ is an isomorphism. Hence $\Gamma \to V'$ is \'etale,
universally injective, and surjective, so an isomorphism by
Morphisms of Spaces,
Lemma \ref{spaces-morphisms-lemma-etale-universally-injective-open}.
Thus $\Gamma$ is the graph of a morphism $a : V' \to V$ as desired.

\medskip\noindent
Finally, we prove that the functor is essentially surjective.
Suppose that $U$ is an object of $X_{spaces, \etale}$.
We have to find an object $V$ of $Y_{spaces, \etale}$
such that $V_X \cong U$. Let $U' \to U$ be a surjective \'etale morphism
such that $U' \cong V'_X$ and $U' \times_U U' \cong V''_X$
for some objects $V'', V'$ of $Y_{spaces, \etale}$.
Then by fully faithfulness of the functor we obtain morphisms
$s, t : V'' \to V'$ with $t_X = \text{pr}_0$ and $s_X = \text{pr}_1$
as morphisms $U' \times_U U' \to U'$. Using that
$(\text{pr}_0, \text{pr}_1) : U' \times_U U' \to U' \times_S U'$
is an \'etale equivalence relation, and that $U' \to V'$ and
$U' \times_U U' \to V''$ are universally injective and surjective
we deduce that
$(t, s) : V'' \to V' \times_S V'$ is an \'etale equivalence relation.
Then the quotient $V = V'/V''$ (see
Spaces, Theorem \ref{spaces-theorem-presentation})
is an algebraic space $V$ over $Y$. There is a morphism
$V' \to V$ such that $V'' = V' \times_V V'$. Thus we obtain a morphism
$V \to Y$ (see
Descent on Spaces, Lemma
\ref{spaces-descent-lemma-fpqc-universal-effective-epimorphisms}).
On base change to $X$ we see that we have a morphism $U' \to V_X$
and a compatible isomorphism $U' \times_{V_X} U' = U' \times_U U'$, which
implies that $V_X \cong U$ (by the lemma just cited once more).

\medskip\noindent
Pick a scheme $W$ and a surjective \'etale morphism $W \to Y$.
Pick a scheme $U'$ and a surjective \'etale morphism $U' \to U \times_X W_X$.
Note that $U'$ and $U' \times_U U'$ are schemes \'etale over $X$ whose
structure morphism to $X$ factors through the scheme $W_X$.
Hence by
\'Etale Cohomology,
Theorem \ref{etale-cohomology-theorem-topological-invariance}
there exist schemes $V', V''$ \'etale over $W$ whose base change to
$W_X$ is isomorphic to respectively $U'$ and $U' \times_U U'$.
This finishes the proof.
\end{proof}

\end{document}
